\subsubsection{Orthogonalität und Winkel}
\label{subsubsec:vector_algebra_orthogonality_angles}

Das Orthogonalitätskriterium aus
Merke~\ref{key:vector_algebra_dot_product_orthogonal_vectors} verwenden wir,
um eine orthogonale Richtung zu bestimmen.

\begin{ExampleBox}{Eine orthogonale Richtung bestimmen}{vector_algebra_find_orthogonal_vector}
  \noindent
  \begin{minipage}[c]{0.57\linewidth}
    Für welches $x\in\R$ sind die beiden Vektoren
    \[
      \vec a=\begin{pmatrix}1\\2\end{pmatrix},
      \qquad
      \vec b=\begin{pmatrix}4\\x\end{pmatrix}
    \]
    senkrecht zueinander?
    \begin{ClozeArea}[20]{9}
      \[
        \begin{array}{@{}rcl@{}}
          \vphantom{\begin{pmatrix}1\\2\end{pmatrix}}\vec a\perp\vec b
          &\Longleftrightarrow& \vec a\cdot\vec b=0\\[0.5em]
          \vphantom{\begin{pmatrix}1\\2\end{pmatrix}}
          &\Longleftrightarrow&
            \begin{pmatrix}1\\2\end{pmatrix}\!\cdot\!
            \begin{pmatrix}4\\x\end{pmatrix}=0\\[0.5em]
          \vphantom{\begin{pmatrix}1\\2\end{pmatrix}}
          &\Longleftrightarrow& 4+2x=0\\[0.5em]
          \vphantom{\begin{pmatrix}1\\2\end{pmatrix}}
          &\Longleftrightarrow& x=-2.
        \end{array}
      \]
    \end{ClozeArea}
  \end{minipage}\hfill
  \begin{minipage}[c]{0.39\linewidth}
    \centering
    \CachedTikzFigure{content/resources/chapters/vector_algebra/tikz/orthogonal_vectors_example.tex}
  \end{minipage}
\end{ExampleBox}

Mit dem Skalarprodukt lässt sich nicht nur die Orthogonalität zweier Vektoren
prüfen, sondern auch ihr Zwischenwinkel berechnen.

\begin{KeyPointBox}{Winkel zwischen zwei Vektoren}{vector_algebra_angle_between_vectors}
  
    Durch Umstellen der geometrischen Definition des Skalarprodukts nach dem
    Zwischenwinkel $\varphi$ erhalten wir
    \par\smallskip\noindent
    \begin{minipage}[c]{0.72\linewidth}
      \[
        \vec a\cdot\vec b=|\vec a|\,|\vec b|\cos(\varphi)
        \quad \Longleftrightarrow \quad
        \varphi=\arccos\!\left(
          \frac{\vec a\cdot\vec b}{|\vec a|\,|\vec b|}
        \right),
        \quad \vec a,\vec b\ne\vec 0.
      \]
  \end{minipage}\hfill
  \begin{minipage}[c]{0.25\linewidth}
    \centering
    \CachedTikzFigure{content/resources/chapters/vector_algebra/tikz/dot_product_geometry.tex}
  \end{minipage}
  \par\smallskip
  Das Skalarprodukt im Zähler kann mithilfe der Koordinatenformel
  ausgewertet werden.
\end{KeyPointBox}

\begin{ExampleBox}{Winkel zweier Vektoren}{vector_algebra_angle_example}
  {\raggedright
    Für die Vektoren$\vec a=\begin{pmatrix}1\\2\end{pmatrix}$ und
    $\vec b=\begin{pmatrix}4\\0\end{pmatrix}$ können wir den
    Zwischenwinkel $\varphi$ bestimmen.\par
  }
  \begin{minipage}[c]{0.69\linewidth}
    \IfClozeTF{%
      \begin{ClozeArea}[24]{13}
      \end{ClozeArea}
    }{%
    \[
      \begin{aligned}
        \varphi
        &=\arccos\!\left(
            \frac{\vec a\cdot\vec b}{|\vec a|\,|\vec b|}
          \right)
          =\arccos\!\left(
            \frac{
              \begin{pmatrix}1\\2\end{pmatrix}\!\cdot\!
              \begin{pmatrix}4\\0\end{pmatrix}}
              {\sqrt{1^2+2^2}\,\sqrt{4^2+0^2}}
          \right)\\[0.8em]
        &=\arccos\!\left(
            \frac{1\cdot4+2\cdot0}
              {\sqrt{1^2+2^2}\,\sqrt{4^2+0^2}}
          \right)
          =\arccos\!\left(
            \frac{4+0}{\sqrt5\,\sqrt{16}}
          \right)\\[0.8em]
        &=\arccos\!\left(\frac{4}{\sqrt5\cdot4}\right)
          =\EquationBranches[\approx]
            {\arccos\!\left(\frac1{\sqrt5}\right)}
            {63{,}4^\circ}
            {\text{(Gradmaß)}}
            {1{,}107}
            {\text{(Bogenmaß)}}
      \end{aligned}
    \]
    }
  \end{minipage}\hfill
  \begin{minipage}[c]{0.29\linewidth}
    \centering
    \CachedTikzFigure{content/resources/chapters/vector_algebra/tikz/angle_between_vectors_example.tex}
  \end{minipage}
\end{ExampleBox}
